Однажды один человек соскочил с трамвая, да так неудачно, что попал под автомобиль. Движение уличное остановилось, и милиционер принялся выяснять, как произошло несчастье. Шофёр долго что"=то объяснял, показывая пальцами на передние колёса автомобиля. Милиционер ощупал эти колёса и записал в книжечку название улицы. Вокруг собралась довольно многочисленная толпа. Какой"=то человек с тусклыми глазами всё время сваливался с тумбы. Какая"=то дама всё оглядывалась на другую даму, а та, в свою очередь, всё оглядывалась на первую даму. Потом толпа разошлась, и уличное движение восстановилось.

Гражданин с тусклыми глазами ещё долго сваливался с тумбы, но наконец и он, отчаявшись, видно, утвердиться на тумбе, лёг просто на тротуар. В это время какой-то человек, нёсший стул, со всего размаха угодил под трамвай. Опять пришёл милиционер, опять собралась толпа и остановилось уличное движение. И гражданин с тусклыми глазами опять начал сваливаться с тумбы.

Ну, а потом опять всё стало хорошо, и даже Иван Семёнович Карпов завернул в столовую.

\begin{flushright}
<\dots{}>
\end{flushright}